% Clase 1 --- Objeto neutro visto de lejos y de cerca (§5).
% (a) Lejos: los versores r_i0 son casi paralelos y las distancias casi iguales,
%     se los saca de factor común y queda sum(Q_i) = 0  =>  F ~ 0.
% (b) Cerca: el 1/r^2 pesa, las cargas próximas dominan y el neto no es nulo.
\begin{center}
\begin{tikzpicture}[scale=1]

  % =================== (a) lejos ===================
  \begin{scope}
    % cúmulo de cargas (neutro)
    \draw[figaux] (0,0) circle (0.95);
    \node[figpos] (p1) at (-0.45,0.42) {};
    \node[figneg] (n1) at (0.38,0.5)   {};
    \node[figpos] (p2) at (0.3,-0.45)  {};
    \node[figneg] (n2) at (-0.35,-0.3) {};
    \node[figpos] (p3) at (-0.1,0.05)  {};
    \node[figneg] (n3) at (0.62,-0.05) {};
    \node[figlbl, above] at (0,1.1) {$\displaystyle\sum_i Q_i = 0$};
    \node[figlblaux, below] at (0,-1.15) {objeto neutro};

    % carga externa, lejos
    \node[figpos] (q0) at (6.6,0) {};
    \node[figlbl, above=4pt] at (q0) {$Q_0$};

    % versores r_i0: casi paralelos vistos desde lejos
    \draw[figvecaux] (p1) -- (6.15,0.14);
    \draw[figvecaux] (p2) -- (6.15,-0.10);
    \draw[figvecaux] (n2) -- (6.15,-0.02);
    \node[figlblaux, above] at (3.5,0.35)
      {$\hat r_{i0}$ casi paralelos, $r_{i0}\simeq r$};

    \node[figlbl, figred] at (3.3,-1.35)
      {$\displaystyle \vec F \simeq \frac{Q_0}{4\pi\varepsilon_0}\,
        \frac{\hat r}{r^2}\sum_i Q_i = 0$};
    \node[figlbl] at (3.3,-2.3) {(a) \textbf{lejos}: la fuerza se cancela};
  \end{scope}

  % =================== (b) cerca ===================
  \begin{scope}[yshift=-5.0cm]
    % objeto polarizado: negativas a la izquierda, positivas a la derecha
    \draw[figaux] (0.7,0) ellipse (1.7 and 0.9);
    \foreach \y in {-0.4,0,0.4} { \node[figneg] at (-0.2,\y) {}; }
    \foreach \y in {-0.4,0,0.4} { \node[figpos] at (1.6,\y) {}; }
    \node[figlblaux, below] at (0.7,-1.1) {objeto neutro, carga mal distribuida};

    % carga externa, pegada al lado positivo
    \node[figpos] (c0) at (3.5,0) {};
    \node[figlbl, above=11pt] at (c0) {$Q_0 > 0$};

    % repulsión de las positivas (cerca): grande
    \draw[figvecres] (c0) -- (5.3,0);
    \node[figlbl, figred, anchor=west] at (5.45,0) {repulsión de las $+$ (cerca)};

    % atracción de las negativas (lejos): chica
    \draw[figvec] (c0) -- (2.75,0);
    \node[figlbl, figblue, anchor=west] at (2.9,-0.8) {atracción de las $-$ (lejos)};

    \node[figlbl] at (3.0,-2.1)
      {(b) \textbf{cerca}: neto repulsivo, $F \sim 1/r^3$};
  \end{scope}

\end{tikzpicture}\\[0.5em]
{\small\itshape Como la fuerza va con $1/r^2$, las cargas más próximas pesan
más: hay tantas atracciones como repulsiones, pero \textbf{no son iguales}. Los
objetos neutros sí generan interacción eléctrica; lo que ocurre es que decae más
rápido que $1/r^2$ ---típicamente como $1/r^3$. La idea de que ``un objeto
neutro no hace fuerza'' vale sólo lejos.}
\end{center}
